\documentclass[10pt,a4paper]{amsart}
\usepackage[margin=19mm]{geometry}
\usepackage{amsmath,amssymb,mathtools}
\usepackage[scr=rsfs,cal=euler]{mathalfa}
\usepackage{xfrac}
\usepackage[unicode,hidelinks,bookmarks=false]{hyperref}
\newcommand{\ie}{i.e.\ }
\newcommand{\cf}{cf.\ }
\newcommand{\resp}{respectively\ }
\newcommand{\Subsection}{\subsection}
\providecommand{\nobreakdash}{\nobreak\hbox{-}}
\setlength{\emergencystretch}{2em}
\multlinegap0pt
\usepackage{bm}
\usepackage{bbm}
\usepackage{dsfont}
\usepackage{xspace}
\usepackage{cancel}
\usepackage{mathtools}
\usepackage{amsthm}
\makeatletter
\@ifundefined{thm}{\newtheorem{thm}{Thm}}{}
\@ifundefined{prop}{\newtheorem{prop}[thm]{Proposition}}{}
\makeatother
\newcommand{\fbl}{\text{FBL}}
\title[Complemented subspaces of Banach lattices]{Complemented subspaces of Banach lattices}
\author{David de Hevia, Pedro Tradacete}
\date{Source: arXiv:2505.24084v2 (CC BY 4.0)}
\begin{document}
\maketitle

\noindent\emph{Source and licence.} Complemented subspaces of Banach lattices. Version arXiv:2505.24084v2 (CC BY 4.0), distributed under the Creative Commons Attribution 4.0 International licence, as recorded in the arXiv licence metadata for this exact version and shown on the abstract page https://arxiv.org/abs/2505.24084v2. Full rights notice: licenses/arxiv-2505.24084-CC-BY-4.0.txt.\par\smallskip

\noindent\textit{Editorial scope.} Counted unit: the Prop whose source label is \texttt{prop:kernel-zero} in the article (source0.tex lines 313--315, source label \texttt{prop:kernel-zero}), delivered together with the article's own complete original proof of it (source0.tex lines 316--337, one \texttt{proof} environment of 22 lines). No mathematics is written, repaired or re-derived: statement and proof are the article's own text line for line. Required context retained uncounted: none. Every definition, display and label of those lines is retained unchanged. Omitted from this package and not redistributed: 1--312, 338--1412. Nothing inside a retained excerpt is omitted or altered: every retained body is delivered exactly as the article's lines are written, so no omission receipt of any kind accompanies this package. Citations: the retained excerpts carry no citation command at all, so no bibliography entry of the article is transcribed and no reference is redistributed. Reference adaptation: none was needed and none was made; every reference inside the delivered unit resolves inside it, so no unresolved reference or citation is produced. Printed slips kept literal: no printed word, symbol, hypothesis or number of the counted statement or of its proof was altered. Layout and class: the article's own class is \texttt{amsart} with packages including inputenc, geometry, color, amssymb, mathrsfs, xy, hyperref, comment, parskip; the wrapper is the lane's pinned \texttt{amsart} editorial wrapper at 10pt on A4, which restates the theorem-like environments the retained text needs and adds the article's own macro definitions that the retained text needs (\texttt{\textbackslash fbl}), with extra wrapper packages (bm, bbm, dsfont, xspace, cancel, mathtools). The delivered numbering is the unit's own. Assets: no retained span contains a figure environment, a graphics-inclusion command, a PDF-page inclusion, a table or a TikZ picture; no image file is copied into this package, the package declares no asset dependency and no image file is redistributed with it. Licence: version arXiv:2505.24084v2 of this article is distributed under the Creative Commons Attribution 4.0 International licence, as recorded in the arXiv licence metadata for this exact version and shown on the abstract page https://arxiv.org/abs/2505.24084v2. A full rights notice is delivered at licenses/arxiv-2505.24084-CC-BY-4.0.txt. Characters: every retained excerpt is pure ASCII, so the delivered engine, XeLaTeX with its default Latin Modern text font, reports no missing character.
\begin{prop}\label{prop:kernel-zero}
    Let $X$ be a Banach lattice. $x^*\in X^*$ is a lattice homomorphism if and only if $\text{ker}(\beta)\subset \text{ker}(\widehat{x^*})$, where $\beta:\fbl[X]\to X$ is the unique lattice homomorphism such that $\beta\circ\delta_X=\text{id}_X$.
\end{prop}

\begin{proof}
Suppose that $x^*:X\to\mathbb{R}$ is a lattice homomorphism. Then, $x^*\circ\beta$ and $\widehat{x^*}$ are lattice homomorphisms on $\fbl[X]$ which clearly agree on $\delta_X(X)$. By uniqueness of the extension, they actually coincide on $\fbl[X]$, so $\widehat{x^*}=x^*\circ\beta$. This shows that $\text{ker}(\beta)\subset\text{ker}(\widehat{x^*})$.

Now, assume that $\text{ker}(\beta)\subset\text{ker}(\widehat{x^*})$. Then the functional $\overline{x^*}:\fbl[X]/\text{ker}(\beta)\to \mathbb{R}$, given by $\overline{x^*}(\bar{f}):=f(x^*)$, is a well-defined lattice homomorphism on $\fbl[X]/\text{ker}(\beta)$. On the other hand, note that the mapping $\bar{\beta}:\fbl[X]/\text{ker}(\beta)\to X$ defined by $\bar{\beta}(\bar{f})=\beta f$ is a bijective lattice homomorphism and it is not difficult to determine its inverse explicitly: the composition $Q\delta_X:X\to\fbl[X]/\text{ker}(\beta)$, where $Q:\fbl[X]\to \fbl[X]/\text{ker}(\beta)$ stands for the canonical quotient defined by $Qf:=\bar{f}$. Indeed, for every $x\in X$, we have
$$
\bar{\beta}Q\delta_X(x)=\bar{\beta}(\overline{\delta_X(x)})=\beta\delta_X(x)=x.
$$
Also, for $f\in\fbl[X]$, we have
$$
Q\delta_X\bar{\beta} (\bar{f})=Q\delta_X\beta (f)=Q(f)=\bar{f},
$$
where the second-to-last identity follows from the fact that $\beta\delta_X=id_X$ and $ker(Q)=ker(\beta)$. Thus, $Q\delta_X$ is also a lattice homomorphism, so $\overline{x^*}\circ Q\delta_X$ is a lattice homomorphism on $X$. But observe that
$$
\overline{x^*}\circ Q\delta_X(x)=\overline{x^*}\bigl(\overline{\delta_X(x)}\bigr)=\delta_X(x)(x^*)=x^*(x),\quad \text{for } x\in X,
$$
so $x^*=\overline{x^*}\circ Q\delta_X$ and this shows that $x^*$ is a lattice homomorphism.
%We will see that it is also norm-preserving. Fix any $\bar{f}\in \fbl[X]/\text{ker}(\beta)$. For every $g\in \text{ker}(\beta)$ we have
%$$
%\|\bar{\beta}(\bar{f})\|=\|\beta(f+g)\|\leq \|f+g\|,
%$$
%so by taking infimum over $g\in \text{ker}(\beta)$ we infer that $\|\bar{\beta}(\bar{f})\|\leq \|\bar{f}\|$. Furthermore, it is clear that the inverse of $\bar{\beta}$ is $Q\delta_X$, where $Q:\fbl[X]\to \fbl[X]/\text{ker}(\beta)$ stands for the canonical quotient defined by $Qf:=\bar{f}$. Indeed, for every $x\in X$, $\bar{\beta}Q\delta_X(x)=\bar{\beta}(\overline{\delta_X(x)})=\beta\delta_X(x)=x$. As $\|Q\delta_X\|\leq 1$, we conclude that $\|\bar{\beta}\bar{f}\|=\|\bar{f}\|$ for all $\bar{f}\in \fbl[X]/\text{ker}(\beta)$. 
\end{proof}

\end{document}
